\documentclass{situationist-preprint}

\title{The Inverted Spectrum and Relationalism}
\msid{MET-17-b}
\msfamily{MET-17}
\msdiscipline{Metaphysics}
\mstype{Companion manuscript}
\mskeywords{inverted spectrum; colour relationalism; phenomenal character; perceiver types; misperception}

\begin{document}
\maketitle

\begin{abstract}
If two perceivers could be functionally alike while the experiences one has of ripe tomatoes are like those the other has of grass, which of them sees the tomato's colour? MET-17-a holds that colours are relational properties of objects, relative to types of perceiver and circumstance. We argue that the inverted spectrum is a question about the relation between colour and phenomenal character, not about the reality of colour, and that relationalism answers it without remainder: inverted perceivers belong to different perceiver types, each sees the colour the tomato has for perceivers of that type, and neither misperceives. Perceiver-independent views must instead posit an error that no fact locates.
\end{abstract}

\section{The case and the question}

Locke considered a person whose ideas of violets and marigolds are reversed relative to ours, and observed that the reversal could never be detected, since the person would call the same things blue and yellow as we do \citep[II.xxxii.15]{Locke1690}. \citet{Shoemaker1982} made the case central to functionalism: if two perceivers can share every functional property while differing in the phenomenal character of their colour experiences, phenomenal character is not functionally definable. Call the two perceivers Norm and Invert. Tomatoes look to Invert the way grass looks to Norm, and conversely, while both sort, name and respond to colours in the same way.

The case raises two questions that are easily run together. The first is whether phenomenal character supervenes on functional organization; that is a question about experience, and nothing here depends on its answer. The second is what colour the tomato is, given that it looks one way to Norm and the other way to Invert. Only the second bears on the reality of colours, and it is the one we address.

\section{The relationalist answer}

Definition~2.1 of MET-17-a says that an object is $C$ for perceivers of type $S$ in circumstances of type $K$ if it is disposed to look $C$ to perceivers of type $S$ in $K$, where \enquote{looks} is phenomenal. The definition individuates colours by phenomenal looks and perceivers by visual constitution, and the following is immediate.

\begin{proposition}\label{p:types}
If Invert is possible, Norm and Invert belong to different perceiver types.
\end{proposition}

\begin{proof}
A perceiver type for colour is a class of perceivers to whom objects are disposed to look alike in the same circumstances. The tomato is disposed to look one way to Norm and another to Invert in the same circumstances, so they are not of the same type, whatever their functional organization.
\end{proof}

\begin{corollary}\label{c:veridical}
Both Norm's and Invert's experiences of the tomato are veridical, and neither misperceives.
\end{corollary}

\begin{proof}
The tomato is red for Norm's type and green for Invert's type, in the phenomenal sense of these words. Each perceiver's experience presents the tomato as having the colour it has for perceivers of that perceiver's type, which is the condition for veridicality in Corollary~3.5 of MET-17-a.
\end{proof}

Public colour words are unaffected. When Invert says \enquote{the tomato is red}, the context fixes, as in ordinary colour talk, the default perceiver type, that of the normal majority in daylight (Section~3.1 of MET-17-a), and his utterance is true. His experience is also veridical, in the sense of Corollary~\ref{c:veridical}. What no one can know is that his type differs from the default, and this ignorance does no harm: it is ignorance of a fact about perceivers, not of a fact about the tomato's colour that someone gets wrong. Real cases come close to the imagined one; \citet{NidaRumelin1996} argued that some subjects with a combination of colour-vision anomalies may have red-green inverted experience, and on the present view such subjects are a further perceiver type, not a class of perceivers in permanent undetectable error.

\section{The comparison}

Perceiver-independent views do worse. If the tomato has a single colour independently of perceivers, then at most one of Norm and Invert sees it as it is, and since they are functionally alike, nothing in either their behavior or the tomato decides which. This is the structure of Lemma~3.2 of MET-17-a: an error with no fact to locate it. Both experiences track the same reflectance type in the same way; only their phenomenal character differs. A view on which colour is fixed by reflectance alone must then say that phenomenal character is irrelevant to colour, which concedes that the experiences of the two perceivers are equally accurate, or that one of them is wrong for no reason. Relationalism avoids the dilemma by letting the phenomenal character of experience fix which relational colour is presented, which is what the inverted spectrum, read as a question about colour, requires. \citet{Cohen2009} develops a relationalism of the same general form.

Nor does the case support eliminativism. The eliminativist infers from the undetectability of inversion that colour experience presents nothing in the world. But each experience presents a property the tomato has, relative to its perceiver type, and the experiences differ because the relational properties presented differ.

\begin{thebibliography}{99}\small
\bibitem[Cohen(2009)]{Cohen2009} Cohen, J. (2009). \emph{The Red and the Real: An Essay on Color Ontology}. Oxford University Press.
\bibitem[Locke(1690)]{Locke1690} Locke, J. (1690). \emph{An Essay Concerning Human Understanding}. London.
\bibitem[Nida-R\"umelin(1996)]{NidaRumelin1996} Nida-R\"umelin, M. (1996). Pseudonormal vision: An actual case of qualia inversion? \emph{Philosophical Studies}, 82(2), 145--157.
\bibitem[Shoemaker(1982)]{Shoemaker1982} Shoemaker, S. (1982). The inverted spectrum. \emph{The Journal of Philosophy}, 79(7), 357--381.
\end{thebibliography}

\end{document}
